% pilot-su2.tex — Chapter on the SU(2) pilot validator: the
% five-criterion inductive closure at finite truncation K=3.
%
% Linked code / Lean:
%   - computations/pilot_su2.py
%   - lean/YMGap/RG/InductiveStep.lean (if present)

\chapter{The \texorpdfstring{$\SUtwo$}{SU(2)} Pilot Validator}
\label{ch:pilot-su2}

\noindent\emph{The five-criterion validator that bundles every
inductive-step closure into a single executable verification at
finite trajectory truncation $K = 3$.}

This chapter assembles the regime-specific certificates of
\Cref{ch:small-field-kp,ch:large-field-kp,ch:overlap-kp} and the
$\sigh$-budget closure of \Cref{ch:sigma-hat-trajectory} into the
manuscript-level pilot validator for $\SUtwo$ at the canonical
operating point $g_0 = 0.30$, $K = 3$.  The validator is a single
predicate that checks five named criteria simultaneously; its
truth at the canonical pilot is the manuscript-quotable evidence
that the inductive-step machinery closes end-to-end at finite
truncation.  The validator is the natural input of
\Cref{ch:uv-trajectory} (which extends to the infinite-$K$ limit)
and of \Cref{ch:mass-gap-theorem} (which feeds the closed
$\sighinf$-positivity into the OS reconstruction).

\section{The five criteria}
\label{sec:five-criteria}

\marginkey{Validator conjuncts}%

The validator checks the following five criteria at each scale
$k \in \{0, 1, \dots, K-1\}$ of the canonical trajectory.

\begin{enumerate}[label=(P\arabic*)]
  \item \textbf{Small-field convergence.}
    $\rhosf(g_k) > \mu_{\mathrm{conn}}$, so the small-field
    Brydges--Federbush cluster expansion converges
    (\Cref{ch:small-field-kp},
    Theorem~\ref{thm:ksf-convergence}).
  \item \textbf{Large-field convergence.}
    $\muplaq \cdot Z_t^{\mathrm{upper}}(g_k) < 1$, so the
    large-field cluster expansion converges
    (\Cref{ch:large-field-kp},
    Theorem~\ref{thm:klf-convergence}).
  \item \textbf{Overlap-cap finiteness.}
    The buffer-overlap summability constant
    $\Kcap(g_k) \le \Ksf + \Klf < \infty$ at every scale
    (\Cref{ch:overlap-kp}, Theorem~\ref{thm:overlap-glue}).
  \item \textbf{Per-step $\eta$ bound.}
    $\etak \le 1.1 \times 10^{-3}$ at every scale on the
    truncated trajectory, where the right-hand side is a
    project-wide ceiling under which the budget closure is
    guaranteed by margin (\Cref{ch:sigma-hat-trajectory},
    Theorem~\ref{thm:eta-coupling}).
  \item \textbf{Aggregate $\sigh$-budget closure.}
    $\sum_{k < K} \etak < \sigma_0$ in $\Qp$
    (\Cref{ch:sigma-hat-trajectory},
    Theorem~\ref{thm:budget-closure}).
\end{enumerate}

\begin{definition}[$\SUtwo$ pilot validator]
\label{def:pilot-validator}
The \emph{$\SUtwo$ pilot validator} at operating point
$(g_0, K)$ is the predicate
\begin{equation}
  \label{eq:validator-predicate}
  \mathrm{Pilot}_{\SUtwo}(g_0, K)
  \;:\Leftrightarrow\;
  \bigwedge_{k = 0}^{K - 1}
    \bigl(\mathrm{P1} \wedge \mathrm{P2} \wedge \mathrm{P3} \wedge \mathrm{P4} \wedge \mathrm{P5}\bigr)
\end{equation}
on the canonical trajectory generated by the one-loop
$\beta$-function from $g_0$ over $K$ steps.
\end{definition}

\section{The canonical pilot closure}
\label{sec:canonical-pilot-closure}

\marginkey{$K = 3$, $g_0 = 0.30$}%

\begin{theorem}[Canonical $\SUtwo$ pilot closure]
\label{thm:canonical-pilot-closure}
$\mathrm{Pilot}_{\SUtwo}(0.30, 3)$ holds.  Specifically, in the
canonical-$\kappa$ convention of \Cref{ch:heat-kernel}:
\begin{align}
  \label{eq:canonical-pilot-values}
  \rhosf(0.30) &\;\approx\; 8.33 \;>\; \mu_{\mathrm{conn}} \approx 3.08
    && \text{(P1, by factor $\approx 2.7$);} \\
  \muplaq \cdot Z_t^{\mathrm{upper}}(0.30) &\;\approx\; 0.0156 \;<\; 1
    && \text{(P2, by factor $\approx 64$);} \\
  \Kcap(g_k) &\;\le\; \Ksf + \Klf \;\approx\; 1.02
    && \text{(P3, finite at every scale);} \\
  \max_k \etak &\;\approx\; 2.72 \times 10^{-4} \;<\; 1.1 \times 10^{-3}
    && \text{(P4, by factor $\approx 4$);} \\
  \sum_{k < 3} \etak &\;\le\; 9 / 10000 \;<\; 1 = \sigma_0
    && \text{(P5, $\sighinf \ge 9991/10000$).}
\end{align}
\end{theorem}

\begin{proof}[Proof sketch]
Each criterion is a constructive $\Qp$-rational check on the
explicit per-scale values; the proof reduces to evaluation of the
closed-form expressions of \Cref{ch:small-field-kp},
\Cref{ch:large-field-kp}, \Cref{ch:overlap-kp}, and
\Cref{ch:sigma-hat-trajectory} at the canonical-trajectory
$\{g_k\}$ values, then comparison against the rational thresholds.
\end{proof}

\begin{infobox}[main][title={Validator output, canonical pilot}]
\centering
\renewcommand{\arraystretch}{1.15}
\begin{tabular}{rcccccc}
\toprule
$k$ & $g_k$    & P1 & P2 & P3 & P4 & P5 \\
\midrule
0   & $0.3000$ & \checkmark & \checkmark & \checkmark & \checkmark & --- \\
1   & $0.3004$ & \checkmark & \checkmark & \checkmark & \checkmark & --- \\
2   & $0.3009$ & \checkmark & \checkmark & \checkmark & \checkmark & --- \\
\midrule
\multicolumn{6}{r}{Aggregate P5:} & \checkmark \\
\bottomrule
\end{tabular}
\vspace{0.2em}\\
Result: $\mathrm{Pilot}_{\SUtwo}(0.30, 3) = \mathsf{True}$,
$\sighinf \ge 9991/10000$.
\end{infobox}

\begin{mainfigure}
\centering
\includegraphics[width=\textwidth]{inductive-closure}
\caption{Five-criterion validator output across the canonical
$\SUtwo$ pilot at $g_0 = 0.30$, $K = 3$.  Each panel
shows one of the conjuncts (P1)--(P5) along the trajectory with the
convergent-regime margin annotated.  All five criteria are
simultaneously satisfied at every scale, giving
$\mathrm{Pilot}_{\SUtwo}(0.30, 3) = \mathsf{True}$ and the
manuscript-headline $\sighinf \ge 9991/10000$ of
\cref{thm:canonical-pilot-closure}.}
\label{fig:inductive-closure}
\end{mainfigure}

\section{The Wilson pilot}
\label{sec:wilson-pilot}

The Wilson convention ($\kappa = 3$ for $\SUtwo$) admits a wider
small-field convergent regime, $g_{\max}^{\mathrm{sf}} \approx 0.99$,
and the pilot validator closes at correspondingly larger
$g_0$ values.

\begin{theorem}[Wilson $\SUtwo$ pilot closure]
\label{thm:wilson-pilot-closure}
$\mathrm{Pilot}_{\SUtwo}(0.50, 3)$ holds.  In the Wilson convention
the validator outputs $\sum_{k < 3} \etak \le 3/10000$, giving
$\sighinf \ge 9997/10000$ — a slightly tighter bound than the
canonical pilot because the Wilson normalization sits further inside
the convergent regime.
\end{theorem}

\section{Failure modes outside the convergent regime}
\label{sec:pilot-failure-modes}

The validator returns \emph{False} cleanly when the trajectory exits
the convergent regime.  Ordered by the coupling $g_0$ at which each
first binds on the canonical trajectory, the failure modes are:

\begin{itemize}
  \item \textbf{Aggregate budget / $\sigh$-retention (P5) --- the
    binding constraint.}
    As $g_0$ increases, the $\sigh$-retention threshold
    $\sighk \ge \tfrac12 \sigma_0$ is the \emph{first} criterion to
    fail, at $g_0 \approx 0.40$, followed immediately by the budget
    closure $\sum_k \etak < \sigma_0$ at $g_0 \approx 0.41$; beyond
    this the cumulative $\eta$-budget exceeds $\sigma_0$ and the
    inductive step no longer closes.  This is the genuine inductive
    failure mode.
  \item \textbf{Large-field / overlap divergence (P2--P4).}
    The large-field and overlap Kotecký--Preiss certificates lose
    convergence at $g_0 \approx 0.42$, just above the budget
    boundary.
  \item \textbf{Small-field divergence (P1).}
    At $g_0 \ge g_{\max}^{\mathrm{sf}}$, $\rhosf \le \mu_{\mathrm{conn}}$
    and the small-field cluster expansion diverges:
    $\Ksf = \infty$, propagating through $\Kcap \to \Dcap \to \etak
    \to \sum \etak = \infty$.  For canonical $\SUtwo$ this is the
    \emph{last} criterion to fail, at $g_0 \approx 0.485$; it is the
    most familiar divergence but, contrary to first appearances, not
    the one that bounds the convergent regime.
\end{itemize}

A fine-resolution scan of the full five-criterion pilot across
$g_0 \in [0.05, 0.498]$ in both heat-kernel conventions is shown in
\cref{fig:g0-phase-boundary}.  The canonical convergent boundary is
$g_0^\ast \approx 0.40$; the Wilson normalization, sitting deeper in
the convergent regime (\cref{thm:wilson-pilot-closure}), passes
across the entire scanned range.  The canonical operating point
$g_0 = 0.30$ used throughout sits well inside this boundary.

\begin{mainfigure}
\centering
\includegraphics[width=\textwidth]{g0-phase-boundary}
\caption{Convergent-regime boundary of the five-criterion pilot as
a function of the initial coupling $g_0$.  Top panel: pilot
pass/fail across $g_0 \in [0.05, 0.498]$ in the canonical and Wilson
conventions; the canonical boundary $g_0^\ast \approx 0.40$ is set
by the $\sigh$-retention threshold (P5), not by the small-field
divergence (P1, which binds only at $g_0 \approx 0.485$).  Bottom
panel: retained $\sighk$ versus $g_0$, with the
$\tfrac12\sigma_0$ retention threshold marked.  The Wilson
convention passes the entire scanned range, consistent with its
wider small-field window $g_{\max}^{\mathrm{sf}} \approx 0.99$.}
\label{fig:g0-phase-boundary}
\end{mainfigure}

The pilot validator is a high-fidelity precondition check for the
continuum-limit construction of \Cref{ch:continuum-measure}:
$\mathrm{Pilot}_{\SUtwo}(g_0, K) = \mathsf{True}$ at every
$K$ is the structural witness that the inductive step closes at
finite truncation, and that the projective limit of the trajectory
admits a positive $\sighinf$.

\section{Code provenance}
\label{sec:pilot-codeanchor}

\begin{itemize}
  \item Python implementation:
    \codepath{computations/pilot\_su2.py}.  Key entry points:
    \codepath{run\_pilot(g\_0, K, t\_choice, sigma\_hat\_initial)}
    returns a \codepath{PilotSU2Result} record carrying the
    per-scale per-criterion booleans, the aggregate validator
    output of \cref{def:pilot-validator}, and the
    $\Qp$-rational $\sighinf$ lower bound.  Tests in
    \codepath{tests/test\_pilot\_su2.py} verify
    \cref{thm:canonical-pilot-closure} and
    \cref{thm:wilson-pilot-closure} exactly.
  \item Visualization:
    \codepath{visualizations/pilot/su2/inductive-closure.svg}
    overlays the five criteria as a single-panel summary across
    the canonical trajectory, with the convergent regime margin
    annotated.
  \item Lean formalization: the canonical-pilot closure is
    Lean-mirrored in \codepath{lean/YMGap/RG/SigmaHatPilot.lean}
    (\codepath{eta\_pilot\_budget\_closes} and related lemmas),
    with axiom-gate \codepath{[propext]} only.
\end{itemize}
